
\textbf{Théorème 5.3.2.} \uline{L'%
application} @@M0@@ (5.3.1) \uline{est un homomorphisme injectif, et l'homomorphisme}
@@M1@@
\uline{induit, par passage au quotient, un homomorphisme inverse à gauche de} @@M2@@ :
@@M3@@
\uline{Si toute section} @@M4@@ \uline{du faisceau constant} @@M5@@ \uline{est bornée (par exemple si} @@M6@@ \uline{est «~quasi-compact~», ou si l'objet final} @@M7@@ \uline{de} @@M8@@ \uline{est somme finie d'objets connexes), alors} @@M9@@ \uline{est un isomorphisme et} @@M10@@ \uline{son inverse.}

\emph{Démonstration} : @@M11@@ est un homomorphisme, car on a pour tout @@M12@@ l'équation
@@M13@@

Admettons pour l'instant que
@@M14@@
i.e. que @@M15@@ soit bien défini. On a alors
@@M16@@
car par 5.2.2 b) on a :
@@M17@@
pour tout @@M18@@. Puis on montre que
@@M19@@
lorsque toute section @@M20@@ de @@M21@@ est bornée. Comme @@M22@@ est le groupe abélien engendré par les éléments @@M23@@, où @@M24@@ est un
